\documentclass[prd,twocolumn,preprintnumbers,showpacs,nofootinbib,superscriptaddress,notitlepage,floatfix]{revtex4-1}
\usepackage{hyperref}
\usepackage{graphicx}
\usepackage{amsmath}
\usepackage{amssymb}
\usepackage{slashed}
\usepackage{amsfonts}
\usepackage{xcolor}
\usepackage{multirow}
\usepackage[caption=false]{subfig}
\usepackage{array}
\usepackage{mwe}


\usepackage[utf8]{inputenc}
\usepackage[normalem]{ulem}

\DeclareMathOperator\Det{Det}
\DeclareMathOperator\re{Re}
\DeclareMathOperator\im{Im}

%  shortcuts for citation related commands
%
\DeclareRobustCommand{\Eq}[1]{Eq.~\eqref{eq:#1}}
\DeclareRobustCommand{\eq}[1]{Eq.~\eqref{eq:#1}}
\DeclareRobustCommand{\Eqs}[2]{Eqs.~\eqref{eq:#1} and \eqref{eq:#2}}
\DeclareRobustCommand{\eqs}[2]{Eqs.~\eqref{eq:#1} and \eqref{eq:#2}}
\DeclareRobustCommand{\eqss}[3]{Eqs.~\eqref{eq:#1}, \eqref{eq:#2} and \eqref{eq:#3}}
\DeclareRobustCommand{\fig}[1]{Fig.~\ref{fig:#1}}
\DeclareRobustCommand{\figs}[2]{Figs.~\ref{fig:#1} and \ref{fig:#2}}
\DeclareRobustCommand{\app}[1]{App.~\ref{app:#1}}
\DeclareRobustCommand{\sec}[1]{Sec.~\ref{sec:#1}}
\DeclareRobustCommand{\secs}[2]{Secs.~\ref{sec:#1} and \ref{sec:#2}}
\DeclareRobustCommand{\tbl}[1]{Table~\ref{tbl:#1}}
\DeclareRobustCommand{\refcite}[1]{Ref.~\cite{#1}}
\DeclareRobustCommand{\refcites}[1]{Refs.~\cite{#1}}

\DeclareRobustCommand{\eq}[1]{Eq.~(\ref{eq:#1})}
\DeclareRobustCommand{\eqs}[2]{Eqs.~(\ref{eq:#1}) and (\ref{eq:#2})}

% Draft Editing Macros
\newcommand{\kl}[1]{{\color{red}{#1}}}
\newcommand{\kll}[1]{{\color{blue}{#1}}}
\newcommand{\MS}{{\overline{\mathrm{MS}}}}
\newcommand{\df}{\mathrm{d}}
\newcommand{\img}{\mathrm{i}}
\newcommand{\eps}{\epsilon}
\newcommand{\cA}{\mathcal{A}}
\newcommand{\cB}{\mathcal{B}}
\newcommand{\cI}{\mathcal{I}}
\newcommand{\cL}{\mathcal{L}}
\newcommand{\cO}{\mathcal{O}}
\newcommand{\cP}{\mathcal{P}}
\newcommand{\cQ}{\mathcal{Q}}
\newcommand{\nn}{\nonumber}
\newcommand{\bn}{{\bar{n}}}
\newcommand{\as}{\alpha_s}
\newcommand{\LQCD}{\Lambda_\mathrm{QCD}}
\newcommand{\UV}{\mathrm{UV}}
\newcommand{\IR}{\mathrm{IR}}
\newcommand{\Tr}{\mathrm{Tr}}
\newcommand{\TMD}{\mathrm{TMD}}
\newcommand{\bt}{\vec b_T}
\newcommand{\kt}{\vec k_T}
\newcommand{\pt}{\vec p_T}
\newcommand{\qt}{\vec q_T}
\newcommand{\ns}{{\rm ns}}
\newcommand{\pb}{{p{\cdot}b}}
\newcommand{\cR}{{\tilde R}}
\newcommand\bets{\begin{table*}}
\newcommand\eets[1]{\label{tb:#1}\end{table*}}

\begin{document}

%%%%%%%%%%%%%%%%%%%%%%%%%%%%%%%%%%%%%%%%%%%%%%%%%%%%%%%%%%%%%%%%%%%%%%
\title{Parton Distributions from Boosted Fields in the Coulomb Gauge}
%%%%%%%%%%%%%%%%%%%%%%%%%%%%%%%%%%%%%%%%%%%%%%%%%%%%%%%%%%%%%%%%%%%%%%

%%%%%%%%%%%%%%%%%%%%%%%%%%%%%%%%%%%%%%%%%%%%%%%%%%%%%%%%%%%%%%%%%%%%%%
%\collaboration{}

\author{Xiang Gao}
%\email{xgao@anl.gov}
\affiliation{Physics Division, Argonne National Laboratory, Lemont, IL 60439, USA}

\author{Wei-Yang Liu}
% \email{wei-yang.liu@stonybrook.edu}
\affiliation{Center for Nuclear Theory, Department of Physics and Astronomy, Stony Brook University, Stony Brook, New York 11794–3800, USA}

\author{Yong Zhao}
% \email{yong.zhao@anl.gov}
\affiliation{Physics Division, Argonne National Laboratory, Lemont, IL 60439, USA}

\begin{abstract}
We propose a new method to calculate parton distribution functions (PDFs) from correlations of boosted quarks and gluons in the Coulomb gauge.
Compared to the widely used quasi-PDFs defined from gauge-invariant Wilson-line operators, such correlations offer advantages including accessibility to larger off-axis momenta, absence of linear power divergence, and enhanced long-range precision. We verify the validity of this method at next-to-leading order in perturbation theory and use it to calculate the pion valence quark PDF on a lattice with spacing $a=0.06$ fm and valence pion mass $m_\pi=300$ MeV.
Our result agrees with that from the gauge-invariant quasi-PDF at similar precision, achieved with only half the computational cost through a large off-axis momentum $|\vec{p}|\sim2.2$ GeV.
This opens the door to a more efficient way to calculate parton physics on the lattice.
\end{abstract}

\maketitle

One of the top quests in nuclear and particle physics nowadays is to understand the 3D internal structure of the proton.
Over the past five decades, high-energy scattering experiments at facilities including SLAC, COMPASS, HERA, Fermilab, LHC, Jefferson Lab and RHIC have provided state-of-the-art measurement of the proton parton distribution functions (PDFs)~\cite{ParticleDataGroup:2022pth}, which are 1D densities of quarks and gluons in their momentum fraction $x$, as well as the 3D and spin-dependent distributions. The future Electron-Ion Collider (EIC) will continue the endeavor with unprecedented precision~\cite{Accardi:2012qut,AbdulKhalek:2021gbh}.

The experimental pursuit of the proton 3D structure has also motivated its first-principles calculation from lattice quantum chromodynamics (QCD), a Euclidean formulation of quantum field theory. However, for a long time such efforts have been hindered by the real-time dependence of light-cone correlations that define the PDFs, which makes them not directly calculable on a Euclidean lattice with imaginary time.
About a decade ago, a breakthrough was made with the proposal of \textit{large-momentum effective theory} (LaMET)~\cite{Ji:2013dva,Ji:2014gla,Ji:2020ect}, which starts from the quasi-PDF (qPDF) defined as Fourier transform of an equal-time correlation at large proton momentum, and relates it to the PDF through power expansion and effective theory matching~\cite{Ji:2020byp}. Over the past years, among the other methods proposed~\cite{Liu:1993cv,Detmold:2005gg,Braun:2007wv,Chambers:2017dov,Radyushkin:2017cyf,Ma:2017pxb}, LaMET has made the most significant progress in the calculation of PDFs and pioneered the studies of 3D partonic structures~\cite{Ji:2020ect,Constantinou:2020hdm,Boussarie:2023izj}.

At the core of LaMET is the simulation of nonlocal bilinear operators which define the qPDFs~\cite{Ji:2013dva}. For example, the quark bilinear is $O_\Gamma(z)\equiv \bar{\psi}(z)\Gamma W(z,0) \psi(0)$, where $\Gamma$ is a Dirac matrix, and $W(z,0)$ is a spacelike Wilson line that connects $0$ to $z^\mu=(0,\vec{z})$ to make $O_\Gamma(z)$ gauge invariant. 
By construction $O_\Gamma(z)$ must approach the light-cone $t+|\vec{z}|=0$ under a Lorentz boost along the $\vec{z}$-direction, which can be achieved on the lattice by simulating a boosted hadron. One major challenge here is to reach large momentum which controls the power accuracy. 
So far the most widely used method to achieve high momenta is the momentum smearing technique~\cite{Bali:2016lva}.
However, to ensure a smooth Wilson line both $\vec{z}$ and the momentum $\vec{p}$ must be along one spatial axis, which leaves out all the off-axis directions that can be used to reach higher momenta~\footnote{A zig-zag Wilson-line operator was studied in Ref.~\cite{Musch:2010ka}, where rotational symmetry was found to be weakly broken on smeared gauge configurations, but it is hard to quantify such an error.}.
Another important issue is the renormalization of $O_\Gamma(z,a)$ under lattice regularization with spacing $a$, as it includes a linear power divergence $\propto \exp(-\delta m(a)|\vec{z}|)$ with $\delta m(a) \!\sim \!1/a$, which originates from the Wilson-line self-energy~\cite{Dorn:1986dt,Maiani:1991az}. In order to calculate the $x$-dependence of PDFs, such a divergence must be subtracted at all $\vec{z}$~\cite{Ji:2020brr}, and a nontrivial matching onto the $\MS$ scheme~\cite{Holligan:2023rex,Zhang:2023bxs} is required to cancel the associated renormalon $\exp(-m_0|\vec{z}|)$ with $m_0\!\sim\! \Lambda_{\rm QCD}$~\cite{Bali:2013pla}, thus eliminating the ${\cal O}(\Lambda_{\rm QCD}/|\vec{p}|)$ correction in the LaMET expansion. 
In practice, this procedure will exponentially amplify the noise in the matrix elements of $O_\Gamma(z,a)$ as $|\vec{z}|$ grows.

In this work we propose to calculate the PDFs from pure quark and gluon correlations in the Coulomb gauge (CG), within the framework of LaMET. The qPDF defined from such a correlation falls into the same universality class~\cite{Hatta:2013gta,Ji:2020ect} as the gauge-invariant (GI) qPDF, since they both approach the PDF under an infinite Lorentz boost. 
Without the Wilson lines, the time for contracting them with the fields is saved in lattice simulations, and one can use larger off-axis momenta. Moreover, the CG correlations are free from linear divergence and the associated renormalon, which greatly simplifies the renormalization as $z$-independent. Therefore, the exponential decay of the bare correlation and its error at large $|\vec{z}|$ is unaffected, allowing better control of the Fourier transform. At last, one can do the momentum smearing~\cite{Izubuchi:2019lyk,Gao:2020ito} in the CG and compute both GI and CG qPDFs simultaneously, for they share the same quark propagators.

In the following, we first introduce the CG qPDF and derive its LaMET matching to the PDF at next-to-leading order (NLO). Next, we calculate the pion valence quark PDF with CG and GI qPDFs on the same lattice ensemble, with the inclusion of a large off-axis momentum $|\vec{p}|\sim2.2$ GeV for the CG qPDF. 
We verify the consistency of both methods in coordinate and momentum spaces, and demonstrate the above mentioned efficiencies of the CG qPDF. 
Finally, we discuss the broader application of CG correlations in calculating parton physics.

\paragraph*{Definition.} The idea of using CG is not new, as it was first proposed within the LaMET framework for calculating the gluon helicity contribution to the proton spin~\cite{Ji:2013fga,Hatta:2013gta,Ji:2014lra,Yang:2016plb}.
The CG quark qPDF is defined as~\footnote{While the CG correlation was examined for the pion wave function in Ref.~\cite{Gupta:1993vp}, it was not adopted as a formal approach there.}
\begin{align}
    \tilde f(x,P^z,\mu) &= P^z\int_{-\infty}^\infty \frac{dz}{2\pi} e^{i x P^z z} \tilde h(z,P^z,\mu)\,,\\
    \tilde h(z,P^z,\mu) &= {1\over 2P^t}\langle P| \bar{\psi}(z)\gamma^t \psi(0)\Big|_{\vec{\nabla}\cdot \vec{A}=0}|P\rangle\,,
\end{align}
where $z^\mu=(0,0,0,z)$, $|P\rangle$ is a hadron state with $P^\mu=(P^t,0,0,P^z)$ normalized to $\langle P|P\rangle = 2P^t\delta^{(3)}(0)$, and $\mu$ is the $\MS$ scale. The GI qPDF follows a similar definition except that the quark correlator is replaced with $O_{\gamma^t}(z)$. The CG condition $\vec{\nabla}\cdot \vec{A}=0$ is fixed so that the quark correlation can have a nonvanishing matrix element. Due to 3D rotational invariance, the correlator and hadron momentum can be oriented to any spatial direction.

Meanwhile, the quark PDF $f(x,\mu)$ is defined as
\begin{align}
    f(x,\mu) &= \int_{-\infty}^\infty \frac{d\lambda}{2\pi} e^{-i\lambda x} h(\lambda,\mu)\,,\\
    h(\lambda,\mu) &= \frac{1}{2P^+} \langle P| \bar{\psi}(\xi^-)W(\xi^-,0) \gamma^+ \psi(0)|P\rangle\,,
\end{align}
where $\lambda=P^+\xi^-$ and $\xi^- = (t- z)/\sqrt{2}$.  Under an infinite Lorentz boost, the CG reduces to the light-cone gauge $A^+=(A^t+A^z)/\sqrt{2}=0$ with a proper boundary condition, where $W(\xi^-,0)=P\exp\big[-ig\int_0^{\xi^-}d\eta^-\ A^+(\eta^-)\big]$ vanishes, so the qPDF becomes equivalent to the PDF.

\paragraph*{Factorization.} According to LaMET~\cite{Ji:2020ect}, when $P^z\gg \Lambda_{\rm QCD}$ the CG qPDF can be perturbatively matched onto the PDF through a factorization formula~\cite{Izubuchi:2018srq},
\begin{align}\label{eq:fact}
	\tilde f(x,P^z,\mu) &= \int {dy\over |y|}\ C\Big({x\over y}, {\mu\over |y|P^z}\Big) f(y,\mu) \nn\\
	&\qquad\qquad + {\cal O}\Big({\Lambda_{\rm QCD}^2\over x^2P_z^2},{\Lambda_{\rm QCD}^2\over(1-x)^2P_z^2}\Big)\,,
\end{align}
where $C$ is the matching coefficient. The factorization is based on effective field theory principles and can be proved using the Feynman diagram analysis for the GI qPDF~\cite{Ma:2014jla,Ji:2020ect}, which will be an object for future study.

By calculating the NLO corrections to the quark CG qPDF and PDF in a free quark state, we find out that their collinear divergences are identical~\cite{Xiong:2013bka}, which confirms \eq{fact} at the same order. The $\MS$ matching coefficient is obtained as a series in the strong coupling $\alpha_s$,
\begin{align}\label{eq:qpdf_match}
	C\big(\xi, {\mu\over p^z}\big) &= \delta(\xi\!-\!1) + {\alpha_sC_F\over 2\pi}C^{(1)}\big(\xi,{\mu\over p^z}\big) + {\cal O}(\alpha_s^2)\,,
\end{align}
where $C_F=4/3$.
At NLO,
\begin{align}
	&C^{(1)}\big(\xi,{\mu\over p^z}\big) = C^{(1)}_{\rm ratio}\big(\xi,{\mu\over p^z}\big)\\
	& + {1\over 2|1-\xi|} +\delta(1-\xi) \left[-{1\over 2}\ln {\mu^2\over 4p_z^2} +{1\over2} - \int_0^1 d\xi' {1\over 1-\xi'}\right]\,,\nn
\end{align}
where
\begin{align}
    &C^{(1)}_{\rm ratio}\big(\xi,{\mu\over p^z}\big) = \left[P_{qq}(\xi)\ln{4p_z^2\over\mu^2} + \xi - 1\right]_{+(1)}^{[0,1]} \\
    & \quad + \Bigg\{P_{qq}(\xi) \Big[\text{sgn}(\xi)\ln|\xi| + \text{sgn}(1\!-\!\xi)\ln|1\!-\!\xi|\Big] \!+\! \text{sgn}(\xi)\nn\\
    &\quad + {3 \xi-1\over \xi-1} \frac{\tan^{-1}\left(\sqrt{1-2 \xi}/|\xi|\right)}{\sqrt{1-2 \xi}} - {3\over 2|1-\xi|} \Bigg\}_{+(1)}^{(-\infty, \infty)}\nn
\end{align}
corresponds to the ratio scheme~\cite{Orginos:2017kos} that satisfies particle number conservation. Here $P_{qq}(\xi)=(1+\xi^2)/(1-\xi)$, and the plus functions are defined on a domain $D$ as
\begin{align}
	\left[g(x)\right]^D_{+(x_0)} &= g(x) - \delta(x-x_0)\int_D dx'\ g(x')\,.
\end{align}
Note that $C^{(1)}_{\rm ratio}$ is analytical at $\xi=1/2$ despite its form.

With a double Fourier transform of \eq{fact}~\cite{Izubuchi:2018srq}, we also derive a short-distance factorization (SDF):
\begin{align}\label{eq:sdf}
	\tilde h(z,\!P^z,\!\mu) &\!=\! \int\! du\ {\cal C}(u,\!z^2\mu^2) h(u \tilde\lambda,\!\mu) \!+\! {\cal O}(z^2\Lambda_{\rm QCD}^2)\,,
\end{align}
where $\tilde\lambda = zP^z$. 
Like \eq{qpdf_match}, the NLO coefficient is
\begin{align}
	{\cal C}^{(1)}(u,z^2\mu^2) &= {\cal C}^{(1)}_{\rm ratio}(u,z^2\mu^2)  \!+\! \delta(1\!-\!u)\big({1\over2}-{\mathbf{L}_z\over2}\big)\,,
\end{align}
where $\mathbf{L}_z=\ln\big(z^2\mu^2e^{2\gamma_E}/ 4\big)$, and
\begin{align}
    &{\cal C}^{(1)}_{\rm ratio}(u,z^2\mu^2)    = \left[-P_{qq}(u)\mathbf{L}_z - {4\ln(1\!-\!u)\over 1\!-\!u} + 1-u\right]_{+(1)}^{[0,1]} \nn\\
    & + \Bigg[ {3 u-1\over u-1} \frac{\tan^{-1}\left(\sqrt{1\!-\!2 u}/|u|\right)}{\sqrt{1-2 u}} \!-\! {3\over |1-u|} \Bigg]_{+(1)}^{(-\infty, \infty)}\,.
\end{align}

In contrast to the matching for GI correlations, where $u$ is limited to $[-1,1]$~\cite{Radyushkin:2017cyf,Ji:2017rah}, the ${\cal C}^{(1)}(u,z^2\mu^2)$ here is nonzero for $u<0$ and $u>1$. As a result, the Mellin moments of ${\cal C}^{(1)}(u,z^2\mu^2)$ are divergent except for the lowest one, indicating that the Wilson coefficients in the operator product expansion of $\tilde h(z,P^z,\mu)$ are functions not only of $z^2$ and $\mu^2$ but also of $P_z^2$. This feature is distinct from the GI case~\cite{Izubuchi:2018srq} and will be further studied.

\paragraph*{Numerical implementation.}
To test the CG qPDF method, we calculate the pion valence quark PDF on a gauge ensemble in 2+1 flavor QCD 
generated by the HotQCD collaboration~\cite{HotQCD:2014kol} with
Highly Improved Staggered Quarks~\cite{Follana:2006rc}, where the lattice spacing $a=0.06$ fm and volume $L_s^3\times L_t=48^3 \times 64$. We use tadpole-improved clover Wilson valence fermions 
on the hypercubic (HYP) smeared~\cite{Hasenfratz:2001hp} gauge background, with a valence pion mass $m_\pi=300$ MeV. To improve the signal of boosted pions at $\vec{p}=(2\pi)/(L_sa)\vec{n}$, we utilize the momentum-smeared~\cite{Bali:2016lva} pion source with optimized quark boost $\vec{k}$~\cite{Izubuchi:2019lyk,Gao:2020ito}.

\begin{table}
\centering
\begin{tabular}{c|c|c|c|c}
\hline
\hline
$|\vec{p}|$ (GeV)&$\vec{n}$ & $\vec{k}$ & $t_s/a$ & (\#ex,\#sl) \\
\hline
 0   &(0,0,0) & (0,0,0) & 8,10,12 & (1, 16) \\
\hline
     && & 8       & (1, 32) \\
 1.72&(0,0,4) & (0,0,3) & 10      & (3, 96) \\
     && & 12      & (8, 256) \\
\hline
     && & 8       & (2, 64) \\
 2.15&(0,0,5) & (0,0,3) & 10      & (4, 128) \\
     && & 12      & (8, 256) \\
 \hline
     & &  & 8       & (1, 32) \\
 2.24&(3,3,3) & (2,2,2) & 10      & (2, 64) \\
     && & 12      & (4, 128) \\
\hline
\hline
\end{tabular}
\caption{Details of lattice setup, where (\#ex,\#sl) are the numbers of exact and sloppy inversions per configuration.}
\label{tbl:setup}
\end{table}

Since lattice simulations involve sampling from statistical ensembles to estimate expectation values, high-momentum modes tend to have larger statistical fluctuations due to their oscillatory nature. Using an off-axis momentum $\vec{n}=(n_x,n_y,n_z)$ one may achieve the same $|\vec{n}|$ with smaller $n_{x,y,z}$ and less statistics. In this study, we employ 109 gauge configurations and perform multiple exact and sloppy Dirac operator inversions on each of them using the All Mode Averaging technique~\cite{Shintani:2014vja}. We use $\vec{n}=(0,0,0)$, two on-axis $\vec{n}=(0,0,n_z)$ with $n_z=4,5$ which correspond to $|\vec{p}|=1.72$ and $2.15$ GeV, and one off-axis $\vec{n}=(3,3,3)$ which corresponds to $|\vec{p}|=2.24$ GeV. Three time separations $t_s/a=8,10,12$ are computed to eliminate the excited-state contamination. Since the quark propagators are the same, we calculate the GI qPDF with 1-step HYP-smeared Wilson lines and the CG qPDF during contraction at no additional cost. More details of the statistics are shown in \tbl{setup}.

For a 4D lattice of spatial volume $V$, we fix QCD in the CG by finding the gauge transformation $\Omega$ of link variables $U_i(t,\vec{x})$ that minimizes the criterion~\cite{Davies:1987vs,Hudspith:2014oja}
\begin{align}
\label{eq:gribov}
  F[U^\Omega]
 & = \frac{1}{9V} \sum_{\vec{x}} \sum_{i=1,2,3} \big[-\re\Tr\ U^\Omega_i(t,\vec{x})\big]\,,
\end{align}
at a precision of $\sim10^{-7}$ on each time slice $t$.
When the observable is not gauge invariant, the above precision and the presence of Gribov copies~\cite{Gribov:1977wm,Singer:1978dk} will affect the lattice result. Despite some proposals for attacking the Gribov problem~\cite{deForcrand:1994mz,Maas:2008ri,Cooper:2015sza,Burgio:2016nad}, there are still no complete solution. Nevertheless, lattice studies of SU(2) Yang-Mills theory show that the Gribov copies only affect the gluon propagator in the far infrared region~\cite{Maas:2017csm}, which implies that they have little impact on the short-range correlations that dominate the PDF at moderate $x$.

Since QCD has been proved to be renormalizable in the CG~\cite{Zwanziger:1998ez,Baulieu:1998kx,Niegawa:2006ey} without linear divergences~\cite{Niegawa:2006hg}, the renormalization of CG correlators is reduced to quark wave function renormalization, which is multiplicative like the GI case~\cite{Ji:2017oey,Ishikawa:2017faj,Green:2017xeu} except for being $z$-independent.

\begin{figure}
    \centering
    \includegraphics[width=0.9\columnwidth]{sdf.pdf}
    \caption{CG and GI ratios. The bands are obtained by matching the PDF fitted from GI ratios at $n_z=4,5$ and $z\in[3a,6a]$.}
    \label{fig:sdf}
\end{figure}

With the bare matrix elements, we first check the consistency between CG and GI correlations at short distance. With a simple paramterization of the PDF, $f_v(x)\propto x^\alpha (1-x)^\beta$, we fit the ratio of GI correlations~\cite{Orginos:2017kos}
\begin{align}\label{eq:ratio}
	{\cal M}(z,P^z,a)&=\frac{\tilde h^{\rm GI}(z,P^z,a)}{\tilde h^{\rm GI}(z,0,a)}\frac{\tilde h^{\rm GI}(0,0,a)}{\tilde h^{\rm GI}(0,P^z,a)}\,,
\end{align}
at $z\in[3a, 6a]$ and $n_z=4,5$, with the NLO SDF formula. Then, we match the fitted PDF to the CG correlations using \eq{sdf} and compare them to the lattice results in \fig{sdf}. After matching, the fitted PDF can describe the CG ratios within $1\sigma$ error, implying that the PDFs calculated from the CG and GI qPDFs must also be consistent at moderate $x$. The slight deviations could come from different ${\cal O}(z^2\Lambda_{\rm QCD}^2)$ corrections that are ignored in the SDF formulas or simply the statistical fluctuations.

Then we continue to the LaMET analysis. First, both bare CG and GI correlations are renormalized in the hybrid scheme~\cite{Ji:2020brr}, where the ratio scheme in \eq{ratio} is used for $z\le z_s$ with $z_s=4a$ and $2\sqrt{3}a$ for on-axis and off-axis momenta, repsectively. At $z>z_s$, the linear divergence in the GI correlation is subtracted with the method in Ref.~\cite{Gao:2021dbh}, and $m_0$ is fitted with the leading-renormalon resummation (LRR) approach in Refs.~\cite{Holligan:2023rex,Zhang:2023bxs}, leaving an overall renormalization to be fixed by a continuity condition at $z=z_s$~\cite{Ji:2020brr}. Meanwhile, the renormalization of CG correlations is simply accomplished with the continuity condition, which leads to nice continuum limit as verified with a finer lattice in \app{Rnm}.
\fig{hbd} compares the hybrid-scheme CG and GI correlations. Both correlations fall close to zero as $z$ increases, but the errors in the GI case are significantly larger due to the exponential enhancement by renormalization. In contrast, the errors in the CG correlation remain small at large $z$. Although at small $z$ the CG correlation has slightly bigger errors, they are likely improvable with better fixed CG condition. Next, we Fourier transform the correlations to obtain the qPDFs. The discrete data are interpolated with a cubic polynomial, whose uncertainty is small compared to the other sources. For the GI correlation, we extrapolate to $z=\infty$ with a physically motivated model $e^{-m|z|}/\tilde \lambda^d$~\cite{Gao:2021dbh}, which mainly affects the small-$x$ region. On the other hand, this extrapolation has a much smaller effect on the CG qPDF since the central value and error of the correlation are both small at large $z$.

\begin{figure}
    \centering
    \includegraphics[width=0.9\columnwidth]{hz}
    \caption{CG and GI correlations in the hybrid scheme at on-axis momentum $2.15$ GeV and off-axis momentum $2.24$ GeV.}
    \label{fig:hbd}
\end{figure}

Subsequently, we match the qPDFs to the PDF. The NLO hybrid-scheme matching coefficient for the GI qPDF is calculated in Ref.~\cite{Ji:2020brr}, and in the CG case it is
\begin{align}
    C^{(1)}(\xi,z_s,p^z,\mu) &= C^{(1)}_{\rm ratio}\big(\xi,{\mu\over p^z}\big)\\
    &  - \left[ {{\rm Si}[(1-\xi)z_sp^z]\over \pi(1-\xi)}-{1\over2 |1-\xi|}\right]_{+(1)}^{(-\infty, \infty)}\,,\nn
\end{align}
where ${\rm Si}(\lambda)=\int_0^\lambda dt\ \sin t/t$. \fig{match} compares the GI qPDF with LRR to the CG qPDF before and after NLO matching. Despite the noticeable difference between the qPDFs, the matched results converge well at $x>0.25$, validating the universality in LaMET~\cite{Hatta:2013gta,Ji:2020ect}.

Finally, we conclude the analysis of CG qPDFs by including the resummation of small-$x$ logarithms through the PDF evolution~\cite{Gao:2021hxl,Su:2022fiu}, while the resummation of large-$x$ logarithms~\cite{Gao:2021hxl,Ji:2023pba} is postponed. \fig{evo} shows the results at on-axis and off-axis momenta $|\vec{p}|=2.15$ and $2.24$ GeV, respectively, which are compared to the recent global fits by \texttt{xFitter20}~\cite{Novikov:2020snp} and JAM21NLO~\cite{Barry:2021osv}. The error includes scale variation, which is estimated by setting $\mu=2\kappa x|\vec{p}|$ with $\kappa=\sqrt{2},1,1/\sqrt{2}$ in the matching and evolving the results to $\mu=2$ GeV at leading-logarithmic (LL) order. The resummation has a huge impact at $x\lesssim 0.2$ where the parton momentum approaches the infrared region. For $x>0.2$, the lattice results agree with the global fits, although they have larger errors. Since the statistics we use is significantly less than that in Ref.~\cite{Gao:2021dbh} for a similar calculation, there is still much room for improvement. More details on the lattice simulation, test of rotational symmetry, renormalization and matching are provided in the Appendix.

\begin{figure}
    \centering
    \includegraphics[width=0.9\columnwidth]{matchingpz5.pdf}
    \caption{PDFs from the qPDFs after NLO matching.}
    \label{fig:match}
\end{figure}

\begin{figure}
    \centering
    \includegraphics[width=0.9\columnwidth]{fvcomparison.pdf}
    \caption{PDFs from CG qPDFs at LL+NLO accuracy, compared to the \texttt{xFitter20}~\cite{Novikov:2020snp} and JAM21NLO~\cite{Barry:2021osv} fits.}
    \label{fig:evo}
\end{figure}

In summary, we have proposed a new method to calculate the PDF from CG correlations within the framework of LaMET. The factorization relation between the CG qPDF and PDF has been verified at NLO. With an exploratory lattice calculation, we demonstrate the equivalence of this method to the GI qPDF and its advantages in achieving larger off-axis momenta, simpler renormalization and more precise long-range correlations at a lower computational cost.

The nice agreement between CG and GI qPDF methods implies a small effect of the Gribov copies, yet further systematic study is still worthwhile.
To improve the precision, we can increase the statistics and pursue higher off-axis momenta.
One practical issue is the large step size along an off-axis direction, such as $\sqrt{3}a$, which adds to the interpolation error. Nevertheless, using the idea of complementarity~\cite{Ji:2022ezo} we can largely overcome it by reconstructing smooth short-range correlations through the SDF of matrix elements at on-axis momenta.
Besides, the evolution and resummations are similar to those for the GI qPDFs, which will be developed in the future for precision calculations.

Finally, the CG correlations can also be used to calculate broader parton physics such as generalized parton distributions and transverse-momentum distributions (TMDs), which are more computationally demanding than the PDFs. In particular, the TMD calculations will benefit significantly from the absence of staple-shaped Wilson lines---whose storage and contractions consume much memory and time---and simplified renormalization~\cite{Constantinou:2019vyb,Ebert:2019tvc,Shanahan:2019zcq,Green:2020xco,Ji:2021uvr,Zhang:2022xuw,Alexandrou:2023ucc}. Since the boosted quarks in a physical gauge capture the correct collinear partonic degrees of freedom, their 3D correlation should be matchable to the physical TMD~\cite{Ebert:2018gzl,Ebert:2019okf,Ji:2019sxk,Ji:2019ewn,Ebert:2022fmh}, which will be studied in a future work.


\begin{acknowledgments}

We thank Jack Holligan, Xiangdong Ji, Swagato Mukherjee and Peter Petreczky for valuable communications.
This material is based upon work supported by the U.S. Department of Energy, Office of Science, Office of Nuclear Physics through Contract No.~DE-AC02-06CH11357 and No.~DE-FG-88ER40388, and within the frameworks of Scientific Discovery through Advanced Computing (SciDAC) award \textit{Fundamental Nuclear Physics at the Exascale and Beyond} and the Quark-Gluon Tomography (QGT) Topical Collaboration, under contract no.~DE-SC0023646. 
YZ is also partial supported by the 2023 Physical Sciences and Engineering (PSE) Early Investigator Named Award program at Argonne National Laboratory.
We gratefully acknowledge the computing resources provided on \textit{Swing}, a high-performance computing cluster operated by the Laboratory Computing Resource Center at Argonne National Laboratory.
This research also used awards of computer time provided by the INCITE program at Argonne Leadership Computing Facility, a DOE Office of Science User Facility operated under Contract No.~DE-AC02-06CH11357. 
The computation of the correlators was carried out with the \texttt{Qlua} software suite~\cite{qlua}, which utilized the multigrid solver in \texttt{QUDA}~\cite{Clark:2009wm,Babich:2011np}.

\end{acknowledgments}

%\newpage

\appendix

\section{Two-point and three-point functions}\label{app:bare}

To determine the bare matrix elements of pion ground state, we first need the two-point functions $C_{2pt}(t_s;\vec{p})$ which will provide energy spectrum created by the pion source and corresponding overlap amplitudes~\cite{Izubuchi:2019lyk}. We utilize the Guassian momentum smeared sources to improve the signal of boosted pion at momentum $\vec{p}=(2\pi)/(L_sa)\vec{n}$. We use $\vec{n}=(0,0,0)$, two on-axis $\vec{n}=(0,0,n_z)$ with $n_z=4,5$ which correspond to $|\vec{p}|=1.72$ and $2.15$ GeV, and one off-axis $\vec{n}=(3,3,3)$ which corresponds to $|\vec{p}|=2.24$ GeV. The optimized quark boost parameters and statistics are shown in \tbl{setup}. In \fig{effmass}, we show the effective mass evaluated from two-point functions as a function of time separation $t_s$. At $t_s\gtrsim 10a$ the effective mass, dominated by the pion ground state, agree with the short colored lines on the right side estimated from the dispersion relation $E=\sqrt{\vec{p}^2+m_\pi^2}$ with $m_\pi$ = 300 MeV. What's more, the signal of $|\vec{p}|=2.24$ GeV case is compatible to the $2.15$ GeV case though the former only takes half of the statistics. This suggests that the off-axis $\vec{n}$ can achieve the same momentum with less computational cost compared to the on-axis ones.

\begin{figure}
    \centering
    \includegraphics[width=0.4\textwidth]{HISQa06_SS_Effmass_pxnpynpzn_cmp.pdf}
    \caption{The effective mass evaluated from two-point functions as a function of $t_s$ are shown. The short colored lines on the right side are estimated from the disperion relation $E=\sqrt{\vec{p}^2+m_\pi^2}$ with $m_\pi$ = 300 MeV.}
    \label{fig:effmass}
\end{figure}

\begin{figure}
    \centering
    \includegraphics[width=0.4\textwidth]{zigzag.pdf}
    \caption{The black arrows are the zig-zag Wilson lines for GI matrix elements with off-axis momentum.}
    \label{fig:zigzag}
\end{figure}

\begin{figure}
    \centering
    \includegraphics[width=0.4\textwidth]{HISQa06_free_ratio_bxp30_px5py0pz0_X0.pdf}
    \includegraphics[width=0.4\textwidth]{HISQa06_free_ratio_bxyzp20_px3py3pz3_XYZ0.pdf}
    \caption{The $C_{3pt}/C_{2pt}$ ratios $R(\tau;\vec{z},\vec{p},t_s)$ at $\vec{z}=\vec{0}$ of $|\vec{p}|$ = 2.15 and 2.24 GeV.}
    \label{fig:ratiofit}
\end{figure}

To extract the the quasi-PDF matrix elements, we need to compute the three point functions $C_{3pt}(\tau,t_s;\vec{p})$. For the case of CG qPDF, we directly do the contraction of the quark propagators without Wilson line, using space separation $\vec{z}$ along the direction $\vec{n}$. As for the case of GI qPDF, we use straight Wilson lines $\vec{z}=(0,0,z_3)$ for on-axis momentum and zig-zag Wilson lines for the off-axis momentum, as shown in \fig{zigzag}. As a result, the distance of a off-axis separation $\vec{z}=\{b,b,b\}$ is $|\vec{z}|=\sqrt{3}b$, while the total length of the Wilson line is $l=3b$. We construct the ratios $R(\tau,\vec{z},\vec{p},t_s)=C_{3pt}(\tau,t_s;\vec{p})/C_{2pt}(t_s;\vec{p})$ to take the advantage of the correlation between two-point and three-point functions. In the $t_s,\tau \rightarrow \infty$ limit, the ratio gives the ground-state matrix elements. In this work, we have calculated three time separation $t_s$ and done a two-state fit~\cite{Izubuchi:2019lyk} for the ground state extrapolation. In \fig{ratiofit}, we show ratios (data points) at $\vec{z}=\vec{0}$ of the two large momenta and the fitted results (colored bands). The black boxes are the ground state matrix elements, where good agreement and similar precision can be observed, though the $|\vec{p}|$ = 2.24 GeV case only used half of the statistics for $|\vec{p}|$ = 2.15 GeV. This is probably due to the smaller momentum modes along each axis.

\section{Bare matrix elements and rotational symmetry}\label{app:bare}

\begin{figure}
    \centering
    \includegraphics[width=0.4\textwidth]{HISQa06_bm_CG.pdf}
    \caption{The bare matrix elements of CG qPDF.}
    \label{fig:bmCG}
\end{figure}

\begin{figure}
    \centering
    \includegraphics[width=0.4\textwidth]{HISQa06_bm_GI.pdf}
    \includegraphics[width=0.4\textwidth]{HISQa06_pn_dm_zs_GI.pdf}
    \caption{The bare (upper panel) and $dm$ subtracted matrix elements (lower panel) of GI qPDF.}
    \label{fig:bmGI}
\end{figure}

In \fig{bmCG}, we show the bare CG qPDF matrix elements as a function of $|\vec{z}|$. It can be seen that the matrix elements from on-axis and off-axis cases overlap with each other, especially at zero momentum with high precision, which implies that the rotational symmetry is well preserved. The bare matrix elements of GI case are shown in the upper panel of \fig{bmGI}. Though the difference of the large-momentum matrix elements is not obvious due to the large errors, there is noticeable deviation for the precise zero-momentum matrix elements. It is evident that zig-zag Wilson line cannot accurately approximate the straight Wilson line. We note that the length of the zig-zag Wilson line is $l=\sqrt{3}|\vec{z}|$. Therefore, in the lower panel of \fig{bmGI} we show the matrix elements after subtracting the linear divergence $e^{-dm\cdot l}$, where $dm$ can be derived from the heavy quark potential ($dm\cdot a=0.1586$)~\cite{Gao:2021dbh}. As one can see, $(e^{dm\cdot |\vec{z}|})^{\sqrt{3}}$ badly overshoots the linear divergence of matrix elements at off-axis $\vec{z}$, which makes their deviation from the on-axis $\vec{z}$ matrix elements even bigger. The reason could be that the HYP smearing distorted the UV physics within a hypercube and the zig-zag Wilson lines contains so many short links. However, smearing is essential to improving the signal of GI qPDF matrix elements, so this obstacle cannot be bypassed. In summary, to use off-axis momenta with reasonable signal and rotational symmetry, the CG qPDF is the better choice.

\section{Multiplicative renormalization of matrix elements in CG}\label{app:Rnm}

\begin{figure}
    \centering
    \includegraphics[width=0.4\textwidth]{HISQa04_pz0_dm_zs_CG.pdf}
    \caption{Comparison of CG ratios $\tilde{h}(\vec{z},\vec{p},a)/\tilde{h}(\vec{z}_s,\vec{p},a)$ at $\vec{p}=0$ and $a=0.04$ and $0.06$ fm. The spatial separation $\vec{z}$ is along the off-axis direction $\vec{n}=(1,1,1)$.}
    \label{fig:CGRnm}
\end{figure}

It is well known that the GI quasi-PDF matrix elements can be multiplicatively renormalized, by removing the linear divergence $\exp(-\delta m(a)|z|)$ and logarithm divergence $Z(a)$~\cite{Ji:2017oey,Green:2017xeu,Ishikawa:2017faj}. In the case of CG quasi-PDF matrix elements, there is no linear divergence but only the logarithm divergences~\cite{Zwanziger:1998ez,Baulieu:1998kx,Niegawa:2006ey,Niegawa:2006hg}. To check the if the CG matrix elements can be multiplicatively renormalized by a constant $Z(a)$, we compute the $\vec{p}=0$ matrix elements on another gauge ensemble in 2+1 flavor QCD generated by the HotQCD collaboration~\cite{HotQCD:2014kol} with Highly Improved Staggered Quarks~\cite{Follana:2006rc}, where the lattice spacing $a=0.04$ fm and volume $L_s^3\times L_t=64^3 \times 64$~\cite{Gao:2020ito}. We use tadpole-improved clover Wilson valence fermions on the HYP-smeared~\cite{Hasenfratz:2001hp} gauge background, with a valence pion mass $m_\pi=300$ MeV. The spatial separation is along the off-axis direction $\vec{n}=(1,1,1)$, and we utilized a total of 12 configurations, each consisting of 1 exact and 16 sloppy inversions.

In \fig{CGRnm}, we show the comparison of CG matrix elements $\tilde{h}(\vec{z},\vec{p},a)$ divided by $\tilde{h}(\vec{z}_s,\vec{p},a)$ for the two lattices with $a=0.04$ and $0.06$ fm. The $z_s$ are set to be same, 0.42 fm. As one can see, the renormalized matrix elements from the two lattices agree with each other very well (except the two points around $|\vec{z}|=0$ with the most serious discretization effect). This shows that the CG matrix elements can be multiplicative renormalized by a constant $Z(a)$ as expected.



\section{Matrix elements in coordinate space}\label{app:sdf}

\begin{figure}
    \centering
    \includegraphics[width=0.4\textwidth]{sdf2.pdf}
    \caption{Comparison of CG and GI ratios with the inclusion of the off-axis momentum $|\vec{p}|=2.24$ GeV.}
    \label{fig:sdf2}
\end{figure}

Following \fig{sdf}, we compare the CG ratio at the off-axis momentum $|\vec{p}|=2.24$ GeV to the band which is obtained by matching the PDF fitted from the GI ratios. Like the $|\vec{p}|=2.15$ GeV case, the lattice ratio at $|\vec{p}|=2.24$ GeV agrees with the band within $1\sigma$ error, which is already implied by the rotational symmetry in \fig{bmCG}.

\begin{figure}
    \centering
    \includegraphics[width=0.4\textwidth]{hzpz4.pdf}
    \caption{Hybrid scheme matrix element}
    \label{fig:hzpz4}
\end{figure}

\fig{hzpz4} compares the hybrid-scheme CG and GI qPDF matrix elements at $P^z=1.72$ GeV, which again shows more precise long-range correlations in the CG case.

\section{Matching}\label{app:match}

\begin{figure}
    \centering
    \includegraphics[width=0.4\textwidth]{matchingpz4.pdf}
    \caption{PDFs from the qPDFs after NLO matching at $P^z=1.72$ GeV.}
    \label{fig:match2}
\end{figure}

\begin{figure}
    \centering
    \includegraphics[width=0.4\textwidth]{fvpz4LL.pdf}
    \includegraphics[width=0.4\textwidth]{fvpz5LL.pdf}
    \includegraphics[width=0.4\textwidth]{fvp333LL.pdf}    
    \caption{PDFs matched from the CG qPDF at NLO and LL+NLO.}
    \label{fig:NLOvsLL}
\end{figure}

\begin{figure}
    \centering
    \includegraphics[width=0.4\textwidth]{fvpz4.pdf}
    \caption{PDFs from CG qPDFs at $P^z=1.72$ GeV, compared to the global fits.}
    \label{fig:comp2}
\end{figure}

Following \fig{match}, we compare the PDFs calculated from the CG and GI qPDFs at  $P^z=1.72$ GeV at NLO in \fig{match2}. Again, despite the considerable differences between the CG and GI qPDFs, the matched PDFs show significantly improved agreement at moderate $x$.

To demonstrate the effect of resumming small-$x$ logarithm or PDF evolution, we compare the PDFs matched from the CG qPDF at NLO and LL+NLO accuracies in \fig{NLOvsLL}. For LL resummation, we use one-loop evolution of $\alpha_s$, which starts at initial value $\alpha_s(\mu=2{\rm GeV})=0.293$. The resummation makes almost no difference to the PDF at $x>0.4$, but becomes more and more significant as $x$ decreases. Eventually, at $2xP^z \sim 0.8$ GeV where $\alpha_s$ becomes of ${\cal O}(1)$, the scale variation uncertainty becomes out of control.

Finally, for completeness we include a comparison of the PDF calculated from the CG qPDF at $P^z=1.72$ GeV to the global fits. Again we find agreement between lattice and phenomenology at moderate to large $x$, which is slightly better than the two larger momenta cases. Since the errors in the current lattice results are huge, the unquantified power corrections, which should be better suppressed at higher momenta, may just be a less important systematic uncertainty here.

\newpage

\bibliography{xpdf}

\end{document}
